\section*{Erdos Problem \#879 --- Round 2}

\subsection*{1) ROUND-2 OBJECTIVE}
I pursue path \textbf{(C)}: identify a fatal formal inconsistency in the Round-1 formalization, give the \emph{minimal corrected statement}, and then prove a genuine new theorem for the corrected formulation.

The reason this is the right direction is that Round-1 quoted the classical upper bound $G(n)<H(n)$, but under the Round-1 definitions this upper bound is already false for small $n$. So before any further progress on the open lower-bound question can be trusted, the formalization itself has to be repaired.

\subsection*{2) Round-1 FOUNDATION USED}
I use the following specific Round-1 results from the exported file \texttt{erdos\_solutions.tex}:
\begin{itemize}
\item Round-1 Lemma 1 for Problem \#879: under the Round-1 formalization, an elementary argument gave only
\[
G(n)\le H(n)+1.
\]
\item Round-1 Lemma 2 for Problem \#879: every prime $p\in(n/2,n]$ is forced in any maximizing admissible set.
\item Round-1 computation table for Problem \#879: exact values of $G(n)$ were computed for several small $n$.
\end{itemize}
I do \emph{not} reprove those results; I only strengthen and correct the setup.

\subsection*{3) NEW INSIGHT / TOOL (ROUND-2)}
The key new insight is that the Round-1 formalization is not merely incomplete: it is \emph{internally inconsistent} with the quoted upper bound.

The minimal repair has two parts.
\begin{enumerate}
\item The admissible set must be taken inside $\{2,3,\dots,n\}$, not $\{1,2,\dots,n\}$, because allowing $1$ always adds an automatic extra $1$ to the optimum.
\item The $H$-term must use only primes in $(\sqrt n,n]$, namely
\[
H_*(n):=\sum_{\substack{p\text{ prime}\\ \sqrt n<p\le n}}p\ +\ n\,\pi(\sqrt n),
\]
not the larger/incorrect Round-1 prime sum $\sum_{p<n}p$.
\end{enumerate}

For this corrected formulation I prove an \emph{exact upper bound}
\[
G_*(n)\le H_*(n),
\]
and an \emph{exact reduction} of the optimization problem to a weighted packing problem over the primes $\le \sqrt n$.

\subsection*{4) ATTACK PLAN (ROUND-2)}
After Round-1, the exact gap was:
\begin{quote}
``Improve the $H$-type upper bound from $H(n)+1$ toward the quoted $H(n)$.''
\end{quote}
The obstruction was that Round-1 counted \emph{all} primes $<n$, which is too crude, and it also allowed $1$, which injects an unavoidable additive error.

So the precise claims to settle now are:
\begin{enumerate}
\item prove that the Round-1 formulation is false as stated by explicit counterexample;
\item identify the minimal corrected formulation compatible with the quoted classical upper bound;
\item prove the corrected upper bound exactly;
\item reformulate the remaining open lower-bound problem in a sharper optimization language.
\end{enumerate}

\subsection*{5) WORK (ROUND-2)}

For clarity I distinguish the \emph{Round-1 quantities}
\[
G_{\mathrm{R1}}(n):=\max_{\substack{S\subseteq\{1,\dots,n\}\\ \gcd(a,b)=1\ \forall a\ne b\in S}}\ \sum_{a\in S}a,
\qquad
H_{\mathrm{R1}}(n):=\sum_{\substack{p\text{ prime}\\ p<n}}p\ +\ n\,\pi(\sqrt n),
\]
from the \emph{corrected quantities}
\[
G_*(n):=\max_{\substack{S\subseteq\{2,\dots,n\}\\ \gcd(a,b)=1\ \forall a\ne b\in S}}\ \sum_{a\in S}a,
\qquad
H_*(n):=\sum_{\substack{p\text{ prime}\\ \sqrt n<p\le n}}p\ +\ n\,\pi(\sqrt n).
\]

\noindent\textbf{Proposition 1 (the Round-1 formalization is false, and both corrections are necessary).}
Under the Round-1 definitions, the quoted upper bound $G(n)<H(n)$ is false.
More precisely:
\begin{enumerate}
\item With the Round-1 definitions exactly as written,
\[
G_{\mathrm{R1}}(7)\ge 20>17=H_{\mathrm{R1}}(7).
\]
\item Even if one excludes $1$ but keeps the Round-1 definition of $H$, the inequality still fails:
\[
\max_{\substack{S\subseteq\{2,\dots,7\}\\ S\text{ admissible}}}\sum_{a\in S}a\ \ge 19>17=H_{\mathrm{R1}}(7).
\]
\item Even if one corrects the prime sum to $H_*$ but still allows $1$, the clean upper bound fails again:
\[
\max_{\substack{S\subseteq\{1,\dots,8\}\\ S\text{ admissible}}}\sum_{a\in S}a\ \ge 24>23=H_*(8).
\]
\end{enumerate}
Hence the minimal consistent repair is: exclude $1$ \emph{and} replace $H_{\mathrm{R1}}$ by $H_*$.

\noindent\emph{Proof.}
(1) Take
\[
S=\{1,3,4,5,7\}\subseteq\{1,\dots,7\}.
\]
This set is admissible because $1$ is coprime to everything and the nontrivial pairwise gcds are
\[
\gcd(3,4)=\gcd(3,5)=\gcd(3,7)=\gcd(4,5)=\gcd(4,7)=\gcd(5,7)=1.
\]
Thus
\[
\sum_{a\in S}a=1+3+4+5+7=20.
\]
On the other hand,
\[
H_{\mathrm{R1}}(7)=\sum_{p<7}p+7\pi(\sqrt7)=(2+3+5)+7\cdot 1=17.
\]
So $20>17$.

(2) Excluding $1$, take instead
\[
S'=\{3,4,5,7\}\subseteq\{2,\dots,7\}.
\]
The same gcd check shows $S'$ is admissible and
\[
\sum_{a\in S'}a=3+4+5+7=19>17=H_{\mathrm{R1}}(7).
\]
So the Round-1 definition of $H$ is itself wrong for the intended upper bound.

(3) Now keep $H_*$ but allow $1$, and take
\[
T=\{1,3,5,7,8\}\subseteq\{1,\dots,8\}.
\]
Again $T$ is admissible because $1$ is harmless and $3,5,7,8$ are pairwise coprime. Hence
\[
\sum_{a\in T}a=1+3+5+7+8=24.
\]
But
\[
H_*(8)=\sum_{\sqrt8<p\le 8}p+8\pi(\sqrt8)=(3+5+7)+8\cdot 1=23.
\]
So $24>23$.
Therefore both corrections are necessary.
\hfill$\square$

\medskip
\noindent\textbf{Lemma 2 (exact corrected upper bound).}
For every $n\ge 2$,
\[
G_*(n)\le H_*(n)=\sum_{\substack{p\text{ prime}\\ \sqrt n<p\le n}}p\ +\ n\,\pi(\sqrt n).
\]

\noindent\emph{Proof.}
Fix an admissible set $S\subseteq\{2,\dots,n\}$.
Define
\[
T:=\{a\in S:\text{$a$ has at least one prime divisor }q\le \sqrt n\},
\qquad
L:=S\setminus T.
\]
Then $S=T\sqcup L$.

First bound $T$.
For each $a\in T$, define a small prime $q(a)\le \sqrt n$ as follows:
\begin{itemize}
\item if $a$ itself is prime, then necessarily $a\le \sqrt n$, and we set $q(a):=a$;
\item if $a$ is composite, let $q(a)$ be its smallest prime divisor.
\end{itemize}
In either case $q(a)$ is a prime with $q(a)\le \sqrt n$ and $q(a)\mid a$.

I claim that the map $a\mapsto q(a)$ is injective on $T$.
Indeed, if $a,b\in T$ are distinct and $q(a)=q(b)=q$, then $q$ divides both $a$ and $b$, so
\[
\gcd(a,b)\ge q>1,
\]
contradicting admissibility of $S$.
Thus the distinct elements of $T$ use distinct small primes, so
\[
|T|\le \pi(\sqrt n).
\]
Since each element of $T$ is at most $n$,
\[
\sum_{a\in T}a\le n\,\pi(\sqrt n).
\]

Now bound $L$.
By definition, no element of $L$ has any prime divisor $\le \sqrt n$.
I claim every element of $L$ is a prime $>\sqrt n$.
Indeed, let $a\in L$. Since $a\ge 2$, it is either prime or composite.
If $a$ were composite, its smallest prime divisor would be at most $\sqrt a\le \sqrt n$, contradicting $a\in L$.
Therefore $a$ is prime, and since it has no prime divisor $\le \sqrt n$, we must have $a>\sqrt n$.
Hence
\[
L\subseteq \{p\text{ prime}:\ \sqrt n<p\le n\},
\]
so
\[
\sum_{a\in L}a\le \sum_{\substack{p\text{ prime}\\ \sqrt n<p\le n}}p.
\]

Finally,
\[
\sum_{a\in S}a=\sum_{a\in T}a+\sum_{a\in L}a
\le n\pi(\sqrt n)+\sum_{\substack{p\text{ prime}\\ \sqrt n<p\le n}}p
=H_*(n).
\]
Since $S$ was arbitrary, $G_*(n)\le H_*(n)$.
\hfill$\square$

\medskip
\noindent\textbf{Lemma 3 (exact reduction to a small-prime packing problem).}
Let $\mathcal T_n$ be the family of sets $T\subseteq\{2,\dots,n\}$ such that:
\begin{itemize}
\item the elements of $T$ are pairwise coprime;
\item every $t\in T$ has at least one prime divisor $\le \sqrt n$.
\end{itemize}
For $T\in\mathcal T_n$, define
\[
U(T):=\{p\text{ prime}:\ \sqrt n<p\le n\text{ and }p\mid t\text{ for some }t\in T\}.
\]
Then
\[
G_*(n)=\max_{T\in\mathcal T_n}\left(\sum_{t\in T}t+\sum_{\substack{p\text{ prime}\\ \sqrt n<p\le n\\ p\notin U(T)}}p\right).
\]
Equivalently,
\[
G_*(n)=\sum_{\substack{p\text{ prime}\\ \sqrt n<p\le n}}p
\ +\ \max_{T\in\mathcal T_n}\left(\sum_{t\in T}t-\sum_{p\in U(T)}p\right).
\]

\noindent\emph{Proof.}
Fix any admissible set $S\subseteq\{2,\dots,n\}$.
Let
\[
T:=\{a\in S:\text{$a$ has a prime divisor }\le \sqrt n\}.
\]
Then $T\in\mathcal T_n$ by construction, because it is a subset of an admissible set.
Every element of $S\setminus T$ has no prime divisor $\le \sqrt n$, so by the argument in Lemma~2 it must be a prime $p$ with $\sqrt n<p\le n$.
Moreover, such a prime $p$ belongs to $S\setminus T$ if and only if $p\notin U(T)$:
\begin{itemize}
\item if $p\in U(T)$, then some $t\in T$ is divisible by $p$, so $p$ cannot also lie in $S$ because then $\gcd(p,t)>1$;
\item if $p\notin U(T)$, then $p$ shares no prime factor with any element of $T$, and it shares no prime factor with any other large prime, so $p$ can be added freely.
\end{itemize}
Therefore for this particular $S$ we have the exact decomposition
\[
\sum_{a\in S}a=\sum_{t\in T}t+\sum_{\substack{p\text{ prime}\\ \sqrt n<p\le n\\ p\notin U(T)}}p.
\]
This proves
\[
G_*(n)\le \max_{T\in\mathcal T_n}\left(\sum_{t\in T}t+\sum_{\substack{p\text{ prime}\\ \sqrt n<p\le n\\ p\notin U(T)}}p\right).
\]

For the reverse inequality, fix any $T\in\mathcal T_n$ and define
\[
S_T:=T\cup \{p\text{ prime}:\ \sqrt n<p\le n,\ p\notin U(T)\}.
\]
I claim $S_T$ is admissible.
Indeed:
\begin{itemize}
\item elements inside $T$ are pairwise coprime by definition of $\mathcal T_n$;
\item distinct added primes are pairwise coprime;
\item if $p$ is one of the added primes and $t\in T$, then $p\notin U(T)$ implies $p\nmid t$, and since $p$ is prime, this is equivalent to $\gcd(p,t)=1$.
\end{itemize}
Hence
\[
\sum_{a\in S_T}a=\sum_{t\in T}t+\sum_{\substack{p\text{ prime}\\ \sqrt n<p\le n\\ p\notin U(T)}}p,
\]
so taking the maximum over $T$ gives the reverse inequality.
Thus the formula is exact.
The equivalent second formula follows by subtracting from the full large-prime sum.
\hfill$\square$

\medskip
\noindent\textbf{Corollary 4 (explicit constructive lower bound).}
For each prime $q\le \sqrt n$, let
\[
\ell_q:=\lfloor \log_q n\rfloor,
\qquad
m_q:=q^{\ell_q}.
\]
Then
\[
G_*(n)\ge \sum_{\substack{p\text{ prime}\\ \sqrt n<p\le n}}p\ +\ \sum_{\substack{q\text{ prime}\\ q\le \sqrt n}} q^{\ell_q}.
\]
Consequently,
\[
0\le H_*(n)-G_*(n)\le \sum_{\substack{q\text{ prime}\\ q\le \sqrt n}}\bigl(n-q^{\ell_q}\bigr).
\]

\noindent\emph{Proof.}
The set
\[
T_0:=\{m_q:\ q\le \sqrt n,\ q\text{ prime}\}
\]
belongs to $\mathcal T_n$: each $m_q\le n$, each $m_q$ has the single prime divisor $q\le \sqrt n$, and distinct prime powers $m_q,m_{q'}$ are pairwise coprime when $q\ne q'$.
Moreover $U(T_0)=\varnothing$, because no $m_q$ has any prime divisor $>\sqrt n$.
Applying Lemma~3 to $T_0$ gives
\[
G_*(n)\ge \sum_{\substack{p\text{ prime}\\ \sqrt n<p\le n}}p + \sum_{q\le \sqrt n} m_q,
\]
as claimed.
Subtracting this from
\[
H_*(n)=\sum_{\substack{p\text{ prime}\\ \sqrt n<p\le n}}p+n\pi(\sqrt n)
\]
yields the stated upper bound on the gap.
\hfill$\square$

\medskip
\noindent\textbf{Interpretation of Lemma 3.}
The corrected open lower-bound question is now equivalent to the following optimization statement:
\begin{quote}
Is it true that there exists a choice of $T\in\mathcal T_n$ such that
\[
\sum_{t\in T}t-\sum_{p\in U(T)}p\ \ge\ n\pi(\sqrt n)-n^{1+o(1)}?
\]
\end{quote}
This is a strictly sharper formulation than in Round-1: instead of treating the problem as a vague ``large primes plus composites'' tradeoff, it becomes an exact weighted packing problem over numbers carrying at least one small prime.

\subsection*{6) ADVERSARIAL VERIFICATION}
\textbf{(a) Boundary checks for the obstruction.}
The counterexamples in Proposition~1 are completely explicit:
\begin{itemize}
\item $n=7$: $\{1,3,4,5,7\}$ has pairwise gcd $1$ and sum $20>17=H_{\mathrm{R1}}(7)$.
\item $n=7$ without $1$: $\{3,4,5,7\}$ has sum $19>17=H_{\mathrm{R1}}(7)$.
\item $n=8$ with corrected $H_*$ but allowing $1$: $\{1,3,5,7,8\}$ has sum $24>23=H_*(8)$.
\end{itemize}
So the diagnosis really does force \emph{both} corrections.

\textbf{(b) Hidden-assumption check in Lemma~2.}
The only nontrivial claim is that an integer $a\le n$ with no prime divisor $\le \sqrt n$ must be a prime $>\sqrt n$.
This is correct because any composite integer has a prime divisor at most its square root.

\textbf{(c) Hidden-assumption check in Lemma~3.}
For a prime $p>\sqrt n$ and an element $t\in T$, the statement ``$p$ is coprime to $t$'' is equivalent to ``$p\nmid t$'' because $p$ is prime. This is exactly why the completion by all primes outside $U(T)$ is valid.

\textbf{(d) Consistency with Round-1.}
Round-1 Lemma~2 (primes in $(n/2,n]$ are forced in maximizing sets) remains valid under the corrected formulation and is fully consistent with Lemma~3: such primes always lie in the added large-prime part unless one deliberately sacrifices one of them inside an element of $T$, which can happen only if the element is exactly that prime itself since $2p>n$.

\textbf{(e) Computational sanity checks.}
I recomputed exact optimal values under the corrected formulation (exclude $1$, use $H_*$). The outputs were:
\begin{verbatim}
Round-1 inconsistency checks:
n=7: old H_R1(7)=17, but {1,3,4,5,7} has sum 20.
n=7: even without 1, {3,4,5,7} has sum 19 > 17.
n=8: corrected H_*(8)=23, but allowing 1 gives {1,3,5,7,8} with sum 24.

Corrected values:
 n   G_*(n)   H_*(n)
 7     19       22
 8     23       23
10     29       32
20     98      112
30    192      209
50    445      511
\end{verbatim}
These checks support the corrected upper bound and directly falsify the Round-1 formalization.

\subsection*{7) FINAL}
\textbf{UNRESOLVED (BUT STRICTLY ADVANCED).}

Round-2 makes strictly substantive progress beyond Round-1:
\begin{enumerate}
\item it identifies and proves a fatal flaw in the Round-1 formalization of Problem \#879;
\item it gives the minimal corrected formulation compatible with the quoted classical upper bound;
\item it proves the corrected exact upper bound $G_*(n)\le H_*(n)$;
\item it gives an exact reduction (Lemma~3) turning the remaining lower-bound question into a concrete weighted packing problem.
\end{enumerate}

What remains open is the true lower-bound side: whether the packing term in Lemma~3 can always achieve
\[
n\pi(\sqrt n)-n^{1+o(1)}.
\]
That is exactly the surviving core difficulty.

\subsection*{8) COMPLETION ESTIMATE}
\noindent\textbf{COMPLETION: 63\%}

\subsection*{9) REFERENCES}
\begin{itemize}
\item Round-1 exported file \texttt{/mnt/data/erdos\_solutions.tex}, Problem \#879 section.
\item Local exact computations performed in this round (dynamic programming over prime-support masks for small $n$).
\end{itemize}
